\documentclass[twoside]{article}

\usepackage{tabulary,graphicx,times,caption,fancyhdr,amsfonts,amssymb,amsbsy,latexsym,amsmath}
\usepackage[utf8]{inputenc}
\pdfoutput=1
\usepackage{url,multirow,morefloats,floatflt,cancel,tfrupee,textcomp,colortbl,xcolor,pifont}
\usepackage[nointegrals]{wasysym}
\usepackage{graphics}
\usepackage{ulem}
\usepackage{multirow}
\usepackage{float}
\usepackage{inputenc}
\usepackage{xspace}
\usepackage{url}
\usepackage{epstopdf}
\usepackage{amsfonts}
\usepackage{amsmath}
\usepackage{amssymb}
\usepackage{extarrows}
\urlstyle{rm}
\graphicspath{ {./images/} }
\allowdisplaybreaks

\makeatletter


\usepackage{ifxetex}
\ifxetex\else
  \usepackage{dblfloatfix}
\fi

\@ifundefined{subparagraph}{
\def\subparagraph{\@startsection{paragraph}{5}{2\parindent}{0ex plus 0.1ex minus 0.1ex}%
{0ex}{\normalfont\small\itshape}}%
}{}
\def\URL#1#2{\@ifundefined{href}{#2}{\href{#1}{#2}}}


\def\UrlOrds{\do\*\do\-\do\~\do\'\do\"\do\-}%
\g@addto@macro{\UrlBreaks}{\UrlOrds}

\makeatother
\def\floatpagefraction{0.8} 
\def\dblfloatpagefraction{0.8}



\usepackage[paperheight=11in,paperwidth=8.3in,margin=2.5cm,headsep=.7cm,top=2.5cm]{geometry}
\usepackage[T1]{fontenc}
\def\floatpagefraction{0.8}
\widowpenalty 10000
\clubpenalty 10000

\renewenvironment{abstract}
	{\trivlist\item[]\leftskip0pt\par\vskip4pt\noindent
  	\textbf{\abstractname}\mbox{\null}\\}
	{\par\noindent\endtrivlist}

\def\keywords#1{\par\medskip\par\noindent\textbf{Keywords}: #1\par}

\linespread{1.13} \date{} \emergencystretch 8pt

\captionsetup[figure]{labelfont=normal,skip=1.4pt,aboveskip=1pc}
\captionsetup[table]{labelfont=normal,skip=1.4pt}

\makeatletter
\def\author#1{\gdef\@author{\hskip-\tabcolsep%
	\parbox{\textwidth}{\raggedright\bfseries#1\\[1pc]}}}
\def\address[#1]#2{\g@addto@macro\@author{\\\hskip-\tabcolsep\parbox{\textwidth}{\raggedright%
	\normalsize\normalfont\textsuperscript{#1}#2}}}
\let\addresslink\textsuperscript
\def\correspondence#1{\g@addto@macro\@author{\\\hskip-\tabcolsep\parbox{\textwidth}{\raggedright%
	\vspace*{10pt}\normalsize\normalfont~\\#1~\\[12pt]}}}
\def\email#1{\g@addto@macro\@author{\\\hskip-\tabcolsep\parbox{\textwidth}{\raggedright%
	\normalsize\normalfont Emails: #1}}}

\def\title#1{\gdef\@title{\vspace*{-30pt}%
	\raggedright\textbf{\@journaltitle}~\\%
  \raggedright\bfseries\ifx\@articleType\@empty\vspace*{20pt}\else%
  \vspace*{20pt}\@articleType\vspace*{20pt}\\\fi#1}}
\let\@journaltitle\@empty \def\journaltitle#1{\gdef\@journaltitle{{\normalfont\itshape#1}}}
\let\@articleType\@empty \def\articletype#1{\gdef\@articleType{{\normalfont\itshape#1}}}

\let\@runningHead\@empty \def\RunningHead#1{\gdef\@runningHead{{\normalfont #1}}}

\usepackage{fancyhdr}
\fancypagestyle{headings}{\renewcommand{\headrulewidth}{0pt}\fancyhf{}
  \fancyhead[R]{\itshape\@runningHead}
  \fancyfoot[C]{\thepage}}
\pagestyle{headings}

\fancypagestyle{plain}{\renewcommand{\headrulewidth}{0pt}%
	\fancyhf{}\fancyhead[R]{Prepared for submission to Hindawi}
  \fancyfoot[C]{\thepage}}
\makeatother

\usepackage[numbers]{natbib}

\setcounter{secnumdepth}{0}
\usepackage{float,xcolor}

\journaltitle{
Advances in High Energy Physics}
\articletype{Research Article} % Research Article/Review Article/Clinical Study

\begin{document}


\title{Strongly coupled fermions in odd dimensions and the running cut-off $\Lambda_d$}


\author{
		Evangelos G. Filothodoros\addresslink{1}}		
% Affiliation
\address[1]{Institute of Theoretical Physics, Aristotle University of Thessaloniki, Thessaloniki, Greece.}

\correspondence{Correspondence should be addressed to 
    	Evangelos G. Filothodoros; efilotho@physics.auth.gr, vagfil79@gmail.com}

% Emails of author




% Running Head


\maketitle 

% Abstract
\begin{abstract}
I study the fermionic $U(N)$ Gross-Neveu model at imaginary chemical potential and finite temperature for odd $d$ dimensions, in the strong coupling regime, by using the gap (saddle point) equation for the fermion condensate of the model. This equation describes the phase transitions from weak to strong coupling regime. I point out that the higher odd dimensional gap equations are linear combinations of the lower dimensional equations in a way that as the dimension of the model increases the lower dimensions are weaker coupled but still in the strong coupling regime. Interestingly, at a specific value of the chemical potential, exactly in the middle of the thermal windows that separate the fermionic from the bosonic (condensed) state of the fermions, I find the mass of the fermion condensate for $d=3,5,7,9$. An anomaly occurs at the $5$ dimensional theory where it is stronger coupled against other theories in higher dimensions and lower energy. The main idea of this work is that the cut-off $\Lambda$ regulator for the UV divergent parts of the fermion mass saddle point equation, plays the role of a physical parameter. This idea is based on the identity of the asymptotic freedom of the Gross-Neveu model as a toy model for QCD.

% Keywords - if any
\keywords{Gross-Neveu; Strong coupling; Cut-off}
\end{abstract}

\section{Introduction}
One of the most extensively used model for the study of chiral symmetry, bosonisation physics \cite{Fradkin:1994tt} via  statistical transmutation \cite{Wilczek:1981du,Polyakov:1988md} and condensed matter physics, is the Gross-Neveu model and its generalizations in the presence of real and imaginary chemical potential. Recently it has been connected (see e.g. \cite{Karch:2016sxi,Murugan:2016zal,Seiberg:2016gmd,Kachru:2016rui,Meng:2020}) to particle-vortex duality e.g. \cite{Peskin:1977kp}. The extension of fermionic theories to finite temperature thermal field theory has also been considered (see  for example \cite{Giombi:2011kc,Aharony:2012ns} and references therein), in the context  of various models that describe matter coupled to non-abelian Chern-Simons fields but also the more recent works about exploring the symmetry-breaking of conformal theories in the large-charge limit like \cite{Alvarez}. At higher dimensions \cite{Genolini, Gaiotto} there are some anomalies (in non-abelian gauge theories like the $5$ dimensional Yang Mills-theories) where we have instantons, associated with a global $U(1)$ symmetry. It seems that the Gross-Neveu model at imaginary chemical potential and finite temperature has an interesting identity at $5$ dimensions ($4+1$ dimensions for the compact dimension of temperature) that creates an anomaly of the original asymptotic freedom, arises from the gap equations of the model.


\section{Fermions coupled to Chern-Simons in a monopole background at $3d$}

When fermions are coupled to a Chern-Simons gauge field in a monopole background, their system exhibits intriguing properties such as the emergence of anyonic statistics, which refers to statistics that are neither fermionic nor bosonic but can be fractional. The study of theories like this finds applications in condensed matter physics, like the FQHE and topological insulators. Instantons \cite{DiVecchia} come into play when one is studying an effective field theory like the Chern-Simons theory that describes the low-energy physics of the FQHE and provide a framework to understand the fractionalized excitations and correspond to nontrivial tunneling processes between different topological sectors. In these theories one may study the transitions between different fractional quantum Hall states, changes in the statistics of excitations and the generation of gaps in the energy spectrum.

At finite temperature, fermions that are coupled to a gauge field $A_\nu$, have a temporal component $A_0$ that acts as an imaginary chemical potential for the $U(1)$ charge (see e.g. \cite{ZinnJustin:2002ru}). 
In \cite{Filothodoros:2016txa} it was argued that for Dirac fermions coupled to an abelian Chern-Simons field at level $k$ in three Euclidean  dimensions\footnote{My notations follow \cite{ZinnJustin:2002ru}.} the presence of a monopole charge is closely related to the presence of an imaginary chemical potential. To understand this connection consider the following fermionic partition function at finite temperature
\begin{align}
\label{DiracCSPF}
Z_{fer}(\beta,k)&=\int [{\cal D}A_\nu][{\cal D}\bar\psi][{\cal D}\psi]\exp{\left[-S_{fer}(\bar\psi,\psi,A_\nu)\right]}\,,\\
\label{Sf}
S_{fer}(\bar\psi,\psi,A_\nu)&=-\int_0^{\beta}\!\!\!dx^0\!\!\int \!\!d^2\bar{x}\left[\bar{\psi}(\slash\!\!\!\partial -i\slash\!\!\!\!A)\psi+i\frac{k}{4\pi}\epsilon_{\nu\lambda\rho}A_\nu\partial_\lambda A_\rho+...\right]\,.
\end{align}
There are also some fermionic self interactions that their presence presented with dots. I expand the CS field around a static (i.e. time independent) monopole configuration $\bar{A}_\nu$ \cite{Fosco:1998cq}
\begin{equation}
\label{Aexpansion}
A_\nu=\bar{A}_\nu+\alpha_\nu\,,\,\,\,\bar{A}_\nu=(0,\bar{A}_1(\bar{x}),\bar{A}_2(\bar{x}))\,,\,\,\, \alpha_\nu=(\alpha_0(x^0),\alpha_1(x^0,\bar{x}),\alpha_2(x^0,\bar{x}))\,,
\end{equation}
normalized as\footnote{For example, one may consider the theory on $S^1\times S^2$.}

\begin{equation}
\label{monopole}
\frac{1}{2\pi}\int d^2x \bar{F}_{12}=1\,,\,\,\,\bar{F}_{\nu\lambda}=\partial_\nu \bar{A}_\lambda-\partial_\lambda \bar{A}_\nu\,
\end{equation}
and $\alpha_\nu$ is a backround gauge field.
Hence, (\ref{Sf}) describes the possibility of monopole configurations in the fermionic theory (or the attachment of $k$ units of monopole charge to the fermions) as
\begin{equation}
\label{Sfexp}
S_{fer}(\bar\psi,\psi,A_\nu) =-\int_0^\beta \!\!\!dx^0\!\!\int \!\!d^2\bar{x}\left[\bar\psi(\slash\!\!\!\partial-i\gamma_i\bar{A}_i-i\gamma_\nu\alpha_\nu)\psi+i\frac{k}{4\pi}\epsilon_{\nu\lambda\rho}\alpha_\nu\partial_\lambda\alpha_{\rho}+..\right]-ik\int_0^\beta \!\!\!dx^0 \alpha_0\,.
\end{equation}
We can perform the path integral over the CS fluctuations projecting into a sector with fixed total monopole charge. To do this, I suppose that there is a mean field approximation in this sector such that the spatial CS fluctuations balance out the magnetic background gauge filed  $\langle \alpha_i\rangle =-\bar{A}_i$ \cite{Barkeshli:2014ida}.\footnote{It would be interesting to further explain whether an appropriate large-$N$ is also necessary for the validity of such an approximation.} It is comperable to how we may think of a reduction in which the background gauge field's  integral along the thermal circle is fixed to be a constant. I then obtain
\begin{align}
Z_{fer}(\beta,k)&=\int [{\cal D}\alpha_0][{\cal D}\bar\psi][{\cal D}\psi]\exp{\left[\int_0^\beta \!\!\!dx^0 \!\!\int \!\!d^2\bar{x}\left[\bar\psi(\slash\!\!\!\partial-i\gamma_0\alpha_0)\psi+..\right]+ik\int_0^\beta \!\!\!dx^0\alpha_0\right]}\nonumber \\
\label{DiracCSPFfin}
&=\int ({\cal D}\theta)e^{ik\theta}Z_{gc,fer}(\beta,-i\theta/\beta)\,,
\end{align}
where  $\theta=\int_0^\beta dx^0\alpha_0(x^0)$, $Z_{gc,fer}(\beta,-i\theta/\beta)$ is the grand canonical partition function for the fermionic theory and I have used standard formulae from \cite{ZinnJustin:2002ru}. We see that the CS level $k$ plays the role of the eigenvalue $Q$ of the $U(1)$ charge operator. 
Here is an important conclusion that the finite temperature partition function of Dirac fermions coupled to abelian CS gauge field at level $k$ in a monopole background, is equivalent to the canonical partition function of the fermions at fixed fermion number $k$.

This discussion shows that the partition function of fermions coupled to abelian CS in a monopole background is intimately related to the respective canonical partition function at fixed total $U(1)$ charge. This charge (instanton number) is carried by instanton configurations of the theory, which in $5d$ theory appear like particles. This picture is analogous to the above example of monopole operators of $3d$ gauge theories \cite{Genolini}. 

These $5d$ theories (like a Yang-Mills theory with a background gauge field) are infrared free and have a singularity pole in the UV requiring a cut-off regulator $\Lambda$. In this paper we will discuss the physical meaning of this cut-off by using the gap equations for the fermionic Gross-Neveu model at imaginary chemical potential and finite temperature (or reducing the theory on a circle). 


\section{The $U(N)$ fermionic Gross-Neveu at imaginary chemical potential and finite temperature in odd dimensions - The gap equations}

In this work I will examine a $d$ odd dimensional fermionic theory like the Gross-Neveu model at imaginary chemical potential (temporal gauge field) and finite temperature and the presence of an anomaly in the $5d$ theory (the Gross-Neveu model is a toy-model for QCD and its asymptotic freedom). It is well known that fermionic Gross-Neveu model exhibit very different patterns of symmetry breakdown at finite temperature $T$ in the absence of a chemical potential $U(N)$. At low temperatures, the Gross-Neveu model exhibits a parity broken phase that vanishes for a critical temperature.
I nevertheless have shown in \cite{Filothodoros:2016txa} that the situation alters in the presence of the imaginary chemical potential and the corresponding phase structure of the model. I have further observed in \cite{Filothodoros:2016txa}  the relevance of the celebrated Bloch-Wigner function \cite{Zagier1} in our calculations of the gap equations and free energies which for $d$-odd can be expressed as finite sums of Nielsen's generalized polylogarithms \cite{Filothodoros:2016txa, Borwein,Kolbig}, while for $d$-even the corresponding expressions are much more complicated. Moreover, I have already noted the relevance of the $1d$ theories to the physics of the the $3d$ model, the relevance of $1d$ and $3d$ theories to the physics of the $5d$ model, the relevance of $1d$, $3d$ and $5d$ theories to the physics of the $7d$ model e.t.c. For these reasons, I will concentrate on studying the generalizations of the 3$d$ model to {\it odd} $d>3$.
The gap equations of the model in arbitrary odd dimensions are:

\subsection{The gap equations}
The GN model in $d$ Euclidean dimensions is described by the generalization of the $3d$ action \cite {Filothodoros:2018, Thesis Filothodoros, Petkou:1998wd, Christiansen:1999uv}. Here I reduce the $d$ theory on the thermal circle, so I get a $d-1$ effective theory which is valid for distances much larger than the radius $\beta=\frac{1}{T}$.

\begin{align}
 S_{fer} = -\int_0^\beta \!\!\!dx^0\int \!\!d^{d-1}\bar{x} \left[\bar{\psi }^{i}(\slash\!\!\!\partial  -i\gamma_0\alpha)\psi ^{i}
+\frac{G_d}{2({\rm Tr}\mathbb I_{d-1})N}\left (\bar{\psi }^{i}\psi ^{i}\right )^{2} +i\alpha NQ_d\right]\,,
\end{align}
with  $Q_d$ the $N$-normalized $d$-dimensional fermionic number density and $i=1,2,..N$ that plays the role of the eigenvalue of the $U(1)$ charge operator. For odd $d$ we take the dimension of the gamma matrices to be ${\rm Tr}{\mathbb I}_{d-1}=2^{\frac{d-1}{2}}$.


The $d$-dimensional  gap equation for the fermionic condensate becomes (by introducing the auxiliary scalar fields $m_*$ and $\alpha_*$):
\begin{align}
\frac{m_*}{G_d}&=\frac{m_*}{\beta}\sum_{n=-\infty}^\infty\int^\Lambda\!\!\frac{d^{d-1} \bar{p}}{(2\pi)^{d-1}}\frac{1}{\bar{p}^2+(\omega_n-\alpha_*)^2+m_*^2}\,
\end{align}



The main issue with the GN model in $d>3$ is that the gap equation has  a finite number of higher order divergent terms as $\Lambda\rightarrow\infty$, which cannot be simply taken care of by the adjustment/renormalization of the single coupling $G_d$ so we have to deal with this issue by using another renormalization method.

I will discuss below in some detail the cases $d=5$, $d=7$ and $d=9$ in order to exhibit some of the general features of the higher dimensional models. Starting with $d=5$ and using the results of \cite{Filothodoros:2016txa, Filothodoros:2018}, we have the form of the gap equation as
\begin{align}
\label{ge5}
m_*\left[-({\cal M}_5\beta^3+D_3(-z_*))-
\frac{1}{2!}\ln^2\!|z_*|\left(D_1(-z_*)-\frac{2}{3\pi}\gamma\right)\right]&=0\,
\end{align}

where 
\begin{equation}
\frac{{\cal M}_5}{(2\pi)^2}=\frac{1}{G_{5,*}}-\frac{1}{G_5}\,,\,\,\,\,\gamma=\Lambda\beta\,
\end{equation}
and $z_*=e^{-i\alpha_*\beta-m_*\beta}$.

$G_5$ is the bare coupling and $G_{5,*}$ is the zero temperature critical coupling. If ${\cal M}_5>0$ then $G_5>G_{5,*}$ and we are in the strong coupling regime. So a non-zero solution of $m_*$ in {(\ref{ge5})} would break parity by giving mass to elementary fermions.
To derive the gap equation I have dropped the infinite number of terms that go as inverse powers of $\Lambda$ and I have noted that the last term in parenthesis resembles the corresponding three-dimensional gap equation. Remarkably we see signs of a partial deconstruction of the higher dimensional models in terms of lowers dimensional quantities. 

As we go to dimensions $d>3$, the crucial issue is the explicit presence of  the cutoff in the gap equation i.e. compare the $5d$ with the $3d$ case. I emphasised that ${\cal M}_5$ is independent of the cutoff $\Lambda$, hence for a given temperature the $5d$ gap equation is a two-parameter equation for $z_*$. This means that there was no unambiguous way to tune the single coupling constant of the theory, namely the parameter ${\cal M}_5$, in order to obtain a cutoff independent result which is equivalent to the nonrenormalizability of the 5-dimensional theory, unless I use another method.
 
In the presence of an imaginary chemical potential the situation became more interesting, since I encountered again nontrivial zeros of $D_3(-z_*)$ on the unit circle. That means that I was able study the critical theory with ${\cal M}_5=0$. A short excursion in Mathematica  yielded two zeroes for $D_3(-z)$ on the unit circle, in order to find the thermal windows of the theory where the system is bosonized. Remarkably their positions were approximated (an interesting and detailed work about these positions are in \cite{Etienne}) to high accuracy in \cite{Filothodoros:2018} by rational multiples of $\pi$ as
\begin{align}
\nonumber D_3(-e^{-i\beta\alpha_*})=Cl_3(\beta\alpha_*\pm\pi)=0\Rightarrow\\ \beta\alpha_*\approx \frac{7\pi}{13}\,{\rm or}\,\beta\alpha_*=\frac{19\pi}{13} \,\,\,({\rm mod}\,2\pi)\,.
\end{align}
Using the periodic properties of the Clausen functions, the relevant results are 
\begin{equation}
Cl_3\left(\frac{6\pi}{13}\right)=Cl_3\left(\frac{20\pi}{13}\right)=0.000362159\,.
\end{equation}
I saw that this patterns generalises to all dimensions.
The charge is zero and the system has been bosonized. However, in contrast with the analogous situation in $d=3$, a nonzero solution for $m_*$ in the critical case ${\cal M}_5=0$ depends on the arbitrary parameter $\gamma$. 

I then briefly remember the seven-dimensional case which shows how my results generalized to higher dimensions \cite{Filothodoros:2018}. The gap equation is
\begin{align}
m_*\left[(-{\cal M}_7\beta^5+D_5(-z_*))+\frac{1}{3!}\ln^2\!|z_*|\left(D_3(-z_*)+\frac{\gamma^3}{45\pi}\right)
+\frac{1}{4!}\ln^4\!|z_*|\left(D_1(-z_*)-\frac{4\gamma}{15\pi}\right)\right]=0\,
\end{align}
where the parameter $\gamma$ has been defined above, 
and 
\begin{equation}
\frac{3{\cal M}_7}{(2\pi)^3}=\frac{1}{G_{7,*}}-\frac{1}{G_7}\,.
\end{equation}

As before, we clearly see in the first equation of $d=7$ case the presence of terms related to the corresponding three- and five-dimensional gap equations.

Moreover as advertised above we see that it was not possible tuning $\gamma$ to remove the constant terms $D_3(-1)$ and $D_1(-1)$ in the expansion of the gap equation near $m_*=0$, and hence to arrange unambiguously for multicritical behaviour for the effective action. This problem clearly persist for all $d>7$. 

Moving on the non zero chemical potential we can look for zeros of the critical gap equation on the unit circle. Again, their positions are remarkably well approximated, better than in $d=5$, by rational multiples of $\pi$ as
\begin{align}
\nonumber D_5(e^{-i\beta\alpha_*})=Cl_5(\beta\alpha_*\pm \pi)=0\Rightarrow \\
 \beta\alpha_*\approx \frac{26\pi}{51}\,{\rm or}\,\frac{76\pi}{51} \,\,\,({\rm mod}\,2\pi)\,.
\end{align}
The relevant result is 
\begin{equation}
Cl_5\left(\frac{25\pi}{51}\right)=Cl_5\left(\frac{77\pi}{51}\right)=0.000129657\,
\end{equation}

The $9d$ gap equation in the same way is:

\begin{align}
m_*[-({\cal M}_9\beta^7+D_7(-z_*))-\frac{\ln^2|z_*|}{4!}(D_5(-z_*)-\frac{4\gamma_1^5}{525\pi})-\\
-\frac{\ln^4|z_*|}{5!}(D_3(-z_*)+\frac{4\gamma_2^3}{315\pi})-\frac{\ln^6|z_*|}{6!}(D_1(-z_*)-\frac{16\gamma_3}{175\pi})]=0
\end{align}

where $\gamma=\Lambda\beta$ the cut-off and

\begin{equation}
\frac{15{\cal M}_9}{(2\pi)^4}=\frac{1}{G_{9,*}}-\frac{1}{G_9}.
\end{equation}

It is interesting as I mentioned in Introduction that the higher odd dimensional gap equation (the $9d$ here) is a linear combination of the lower dimensional equations ($3d$, $5d$ and $7d$). One may write:

\begin{equation}
m_*(-g_9-\frac{\ln^2|z_*|}{4!}g_7
-\frac{\ln^4|z_*|}{5!}g_5-\frac{\ln^6|z_*|}{6!}g_3)=0
\end{equation}
and for $m_*\not=0$
\begin{equation}
\label{gaplinear9}
g_9=-\frac{\ln^2|z_*|}{4!}g_7
-\frac{\ln^4|z_*|}{5!}g_5-\frac{\ln^6|z_*|}{6!}g_3
\end{equation}
where $g_d$ the gap equation for $d=3,5,7,9$.


Again for non-zero values of the chemical potential we have the thermal window's opening and closing values of $\beta\alpha_*$

\begin{align}
\nonumber D_7(e^{-i\beta\alpha_*})=Cl_7(\beta\alpha_*\pm \pi)=0\Rightarrow \\
 \beta\alpha_*\approx \frac{103\pi}{205}\,{\rm or}\,\frac{307\pi}{205} \,\,\,({\rm mod}\,2\pi)\,.
\end{align}
The relevant result is 
\begin{equation}
Cl_7\left(\frac{102\pi}{205}\right)=Cl_7\left(\frac{308\pi}{205}\right)=-0.000101475\,
\end{equation}

The basic features discussed above do not change as we move to higher dimensions. We continue to see the partial deconstruction of the $d$-dimensional gap equations in terms of lower-dimensional pieces. Namely, the $d$ dimensional gap equation contains the $d-2, d-4,...,5,3$-dimensional gap equations.

The thermal windows are:\\
	
	{\centering
		\begin{tabular}{|c|c|c|}
			\hline 
			\multicolumn{3}{|c|}{\textbf{Table 1. Thermal Windows for the Gross-Neveu model}} \\ 
			\hline 
			Dimensions & Closing T & Opening T \\ 
                        \hline
                        3 & $\frac{3\alpha_*}{4\pi}$ & $\frac{3\alpha_*}{2\pi}$\\
			\hline 
			5 & $\frac{13\alpha_*}{19\pi}$ & $\frac{13\alpha_*}{7\pi}$ \\ 
			\hline 
			7 & $\frac{51\alpha_*}{76\pi}$ & $\frac{51\alpha_*}{26\pi}$ \\ 
			\hline 
			9 & $\frac{205\alpha_*}{307\pi}$ & $\frac{205\alpha_*}{103\pi}$ \\ 
			\hline 
			11 & $\frac{819\alpha_*}{1228\pi}$ & $\frac{819\alpha_*}{410\pi}$ \\ 
			\hline 
		\end{tabular}\par} 

Let's focus on the case of the $3d$ theory thermal window. The borders for $\alpha_*$ are $2\pi T/3$ and $4\pi T/3$. These are the points where the $D_2(-z*)$ \cite{Zagier1, Zagier2} takes its maximum value (imaginary part) on the unit circle. On the unit circle $D_2(-z*)=Cl_2(\pi\pm\beta\alpha_*)$. At the middle point of the thermal windows $D_{even}(-z*)$ is zero and we are at the full bosonisation points.

\section{Strongly coupled fermions and anomaly of the $5d$ gap equation}
 To calculate the condensate gap-equation of the $U(N)$ Gross-Neveu model in arbitrary odd dimensions and in the presence of imaginary chemical potential $\mu=-i\alpha$ I use the Euclidean action from \cite{Filothodoros:2018}.
 \begin{equation}
 	S_{fer} = -\int_0^\beta \!\!\!dx^0\int \!\!d^{d-1}\bar{x} \left[\bar{\psi }^{i}(\slash\!\!\!\partial  -i\gamma_0\alpha)\psi ^{i}+\frac{G_d}{2({\rm Tr}\mathbb I_{d-1})N}\left (\bar{\psi }^{i}\psi ^{i}\right )^{2} +i\alpha NQ_d\right]\,,\,\,\,i=1,2,..N.
 \end{equation}
where $Q_d$ is the eigenvalue of the $N$-normalized fermion number density operator $\hat{Q}_d=\psi^{i\dagger}\psi^i/N$ in $d$ odd dimensions and it comes from the Lagrangian of the model that possesses a $U(1)$ global symmetry. Introducing an auxiliary scalar field $m$ the canonical partition function is given by
 \begin{equation}
 	S_{fer,eff}=iQ_d\int_0^\beta \!\!\!dx^0\!\!\int\!\!d^{d-1}\bar{x}\,\alpha-\frac{{\rm Tr}\mathbb I_{d-1}}{2G_d}\int_0^\beta \!\!\!dx^0\int d^{d-1}\bar{x} m^2+\rm Tr\ln(\slash\!\!\!\partial-i\gamma_0\alpha+m)_\beta
 \end{equation}
 
 To evaluate the condensate gap equation I look for constant saddle points $\alpha_*$ and $m_*$. At large-$N$ we have the gap equation
 \begin{equation}
 	\frac{\partial}{\partial m}S_{fer,eff}\Biggl|_{m_*,\alpha_*}\!\!\!\!=0\,\,\,\Rightarrow\,\,\, -\frac{m_*}{G_d}+\frac{m_*}{\beta}\sum_{n=-\infty}^\infty\int^\Lambda\!\!\frac{d^{d-1} \bar{p}}{(2\pi)^{d-1}}\frac{1}{\bar{p}^2+(\omega_n-\alpha_*)^2+m_*^2}=0
 \end{equation}
 where  the fermionic Matsubara sums are over the discrete frequencies $\omega_n=(2n+1)\pi/\beta$. The divergent integrals are regulated by the cutoff $\Lambda$. 
 The main issue with the GN model in $d>3$ is that the gap equation has  a finite number of higher order divergent terms as $\Lambda\rightarrow\infty$, which cannot be simply taken care of by the adjustment/renormalization of the single coupling $G_d$. In order to calculate the above integral of the general $d$ gap equation I use a standard inversion formula for the hypergeometric function which allows me to obtain the cut off $\Lambda$ part 
\begin{align}
\label{Cutoff}
\int^\Lambda\!\!\!\frac{d^d p}{(2\pi)^d}\frac{1}{p^2+m*^2} &= \frac{1}{G_{d,*}}-\frac{\Lambda^{d-2}}{d-2}\frac{S_d}{(2\pi)^d}{}_2F_1\left(1,\frac{d}{2}-1;\frac{d}{2};-\frac{\Lambda^2}{m*^2}\right)\nonumber\\
&\hspace{-2cm}=\frac{1}{G_{d,*}}-\frac{S_d}{(2\pi)^d}\left[\Gamma\left(\frac{d}{2}\right)\Gamma\left(2-\frac{d}{2}\right)\frac{m*^{d-2}}{d-2}+\frac{m^2_*\Lambda^{d-4}}{d-4}{}_2F_1\left(1,2-\frac{d}{2};3-\frac{d}{2};-\frac{m*^2}{\Lambda^2}\right)\right]
\end{align}
where $S_d=2\pi^{d/2}/\Gamma(d/2)$. This way we  see that for odd $d$ there are a finite number of divergent terms as $\Lambda\rightarrow\infty$. To avoid these divergent parts I have to try a kind of regularization (similar to \cite{Filothodoros Lattices}) where the gap equation is $D_{odd}(-z*)=0$ at the highest odd dimension and at the critical point of the corresponding field theory. The lower dimensions have their own gap equations inside the gap equation of the $d$ model. A new approach gives to the cut-off $\Lambda$ different values depending on the dimension of the gap equation. It is like we change the bare cut-off $\Lambda$ with a running cut-off $\Lambda_d$ that varies as a function of energy. It takes into account the renormalization group flow of the theory, which describes how the theory's parameters evolve as the energy scale changes, instead of the bare cut-off which is a fixed regularization scale introduced at the beginning of a calculation to handle UV divergences. Let's see for example ({\ref{Cutoff}}) at $d=7$. The last hypergeometric term gives us the sum:

\begin{equation}
\label{Parts}
\frac{m_*^2\Lambda_d^3}{3}-m_*^4\Lambda_d-\frac{m_*^6}{\Lambda_d}
\end{equation}

For large energies (large running cut-off) the last term of ({\ref{Parts}}) goes to zero so, we may always regulate the theory of higher dimension in higher energy in order to eliminate the $\frac{1}{\Lambda_d}$ part. On the other hand the other terms are infrared free, so we have to regularize the theory. As the dimension decreases the running cut-off (energy) increases and the theory is at its weak coupling. On the other hand as the dimension increases the running cut-off decreases and the fermions are more stronger coupled according to asymptotic freedom.

The critical gap equation $({\cal M}_d=0)$ in odd dimensions that determines the condensate mass is (I assume that we are at $T=0$ where we substitute the bare coupling with its renormalized version):

\begin{equation}
-{\cal M}_d\beta^{d-2}+D_{d-2}(-z_*)=0\rightarrow D_{d-2}(-z_*)=0
\end{equation}
for $d=3,7,11,...$

and 

\begin{equation}
{\cal M}_d\beta^{d-2}+D_{d-2}(-z_*)=0\rightarrow D_{d-2}(-z_*)=0
\end{equation}
for $d=5,9,13,...$. 
 
 \begin{itemize}
 	\item \underline{3d gap equation}
 \end{itemize}
 
 The gap equation in 3 dimensions of the Gross-Neveu model at imaginary chemical potential at the critical point ${\cal M}_3=0$ or $G_3=G_{3,*}$ and for $\alpha_*=\pi/\beta$, $m_*\not=0$ turns to:
 
 $m_*\left[-{\cal M}_3\beta+D_1(e^{-m_*\beta})\right]=0\,\rightarrow m_* D_1(e^{-m_*\beta})=0\rightarrow D_1(e^{-m_*\beta})=0\rightarrow \ln(1-e^{-m_*\beta})-\ln\frac{|e^{-m_*\beta}|}{2}=0\rightarrow m_*\beta=\ln\phi^2$ \cite {Filothodoros:2016txa,Sachdev:1993pr}. Also $-z_*=e^{-m_*\beta}$ at $\alpha_*=\pi/\beta$ and $\beta=\frac{1}{T}$. 
 
 
 \begin{itemize}
 	\item \underline{5d gap equation}
 \end{itemize}
 
 The gap equation in 5 dimensions of the Gross-Neveu model at imaginary chemical potential for $\alpha_*=\pi/\beta$, $m_*\not=0$ turns to:
 
\begin{equation}
m_*[-({\cal M}_5\beta^3+D_3(e^{-m_*\beta}))-\frac{\ln^2|z_*|}{2!}(D_1(e^{-m_*\beta})-\frac{2\gamma}{3\pi})]=0
\end{equation}
where $\gamma=\Lambda\beta$ the cut-off. Obviously the gap equation includes the $D_3$ part for the $5d$ theory and the $3d$ gap equation with the $D_1$ part. If we set $\gamma=\Lambda_3\beta$, put the $5d$ theory at the critical point ${\cal M}_5=0$ or $G_5=G_{5,*}$ and assume that $D_1(e^{-m_*\beta})-\frac{2\Lambda_3\beta}{3\pi}=0$, we find from gap equation $D_3(e^{-m_*\beta})=0$: 
\begin{equation}
\boxed
{m_*\beta=2.03185}
\end{equation}
and from the $3d$ equation that $\Lambda_3\beta=4.12525$. 
 
We see that $D_1(e^{-m_*\beta})-\frac{2\Lambda_3\beta}{3\pi}=0$ is the $3d$ gap equation so we may change $\frac{2\Lambda_3\beta}{3\pi}$ with ${\cal M}_3^{(5)}\beta$ which is the mass scale that separates weak from strong coupling of the $3d$ theory included in the $5d$ theory. 
 
 \begin{itemize}
 	\item \underline{7d gap equation}
 \end{itemize}
 
 The gap equation in 7 dimensions of the Gross-Neveu model at imaginary chemical potential and for $\alpha_*=\pi/\beta$, $m_*\not=0$ turns to:
 
\begin{equation}
m_*[(-{\cal M}_7\beta^5+D_5(e^{-m_*\beta}))+\frac{\ln^2|z_*|}{3!}(D_3(e^{-m_*\beta})+\frac{\gamma_1^3}{45\pi})+\frac{\ln^4|z_*|}{4!}(D_1(e^{-m_*\beta})-\frac{4\gamma_2}{15\pi})]=0
\end{equation}
where $\gamma=\Lambda\beta$ the cut-off. Again the gap equation includes the $D_5$ part for the $7d$ theory and the $5d$ and $3d$ equations with $D_3$ and $D_1$ parts respectively. If we set $\gamma_2=\Lambda_3\beta$ and $\gamma_1=\Lambda_5\beta$ and put the $7d$ theory at the critical point ${\cal M}_7=0$ or $G_7=G_{7,*}$, we find from gap equation $D_5(e^{-m_*\beta})=0$:
\begin{equation}
\boxed 
{m_*\beta=2.89218}
\end{equation}
and that $\Lambda_5\beta=6.0403$ from $D_3(e^{-m_*\beta})+\frac{(\Lambda_5\beta)^3}{45\pi}=0$ and $\Lambda_3\beta=16.3642$ from $D_1(e^{-m_*\beta})-\frac{4\Lambda_3\beta}{15\pi}=0$.

We will also change (like the $5d$ case) $\frac{4\Lambda_3\beta}{15\pi}$ with ${\cal M}_3^{(7)}\beta$ and $\frac{(\Lambda_5\beta)^3}{45\pi}$ with ${\cal M}_5^{(7)}\beta^3$, which are the mass scales that separate weak from strong coupling of the $3d$ and $5d$ respectively, included in the $7d$ theory. 
 
 \begin{itemize}
 	\item \underline{9d gap equation}
 \end{itemize}
 
The gap equation in 9 dimensions of the Gross-Neveu model at imaginary chemical potential for $\alpha_*=\pi/\beta$, $m_*\not=0$ turns to:

\begin{align}
m_*[-({\cal M}_9\beta^7+D_7(e^{-m_*\beta}))-\frac{\ln^2|z_*|}{4!}(D_5(e^{-m_*\beta})-\frac{4\gamma_1^5}{525\pi})-\\\nonumber
-\frac{\ln^4|z_*|}{5!}(D_3(e^{-m_*\beta})
+\frac{4\gamma_2^3}{315\pi})
-\frac{\ln^6|z_*|}{6!}(D_1(e^{-m_*\beta})-\frac{16\gamma_3}{175\pi})]=0
\end{align}
where $\gamma=\Lambda\beta$ the cut-off. If we set $\gamma_1=\Lambda_7\beta$, $\gamma_2=\Lambda_5\beta$ and $\gamma_3=\Lambda_3\beta$ and put the $9d$ theory at the critical point ${\cal M}_9=0$ or $G_9=G_{9,*}$ we find from $D_7(e^{-m_*\beta})=0$:

\begin{equation}
\boxed
{m_*\beta=3.68896}
\end{equation}
and that $\Lambda_7\beta=3.8873$ from $D_5(e^{-m_*\beta})-\frac{4(\Lambda_7\beta)^5}{525\pi}=0$, $\Lambda_5\beta=9.8758$ from $D_3(e^{-m_*\beta})+\frac{4(\Lambda_5\beta)^3}{315\pi}=0$ and $\Lambda_3\beta=65.0883$ from $D_1(e^{-m_*\beta})-\frac{16\Lambda_3\beta}{175\pi}=0$.

We will also change (like the $5d$ and $7d$ cases) $\frac{16\Lambda_3\beta}{175\pi}$ with ${\cal M}_3^{(9)}\beta$, $\frac{4(\Lambda_5\beta)^3}{315\pi}$ with ${\cal M}_5^{(9)}\beta^3$ and $\frac{4(\Lambda_7\beta)^5}{525\pi}$ with ${\cal M}_7^{(9)}\beta^5$, which are the mass scales that separate weak from strong coupling of the $3d$, $5d$ and $7d$ respectively, included in the $9d$ theory. All the values of the mass scales were calculated and listed in table 3 as parts of the odd dimensional gap equations we study. For example as we see in table 3, ${\cal M}_3^{(5)}\beta=0.8754$, ${\cal M}_5^{(7)}\beta^3=1.5589$, ${\cal M}_7^{(9)}\beta^5=2.1527$, etc.
 
The higher dimensions gap equations contain the lower dimension equations giving us an idea of what happens to the strong coupling in lower dimensions when in the upper dimension we are at the critical point ${\cal M}_{odd}=0$ (the lower dimensions are weaker coupled but still in  the strong coupling regime, since $G_{odd}>G_{odd,_*})$. The overall picture inside the strong coupling regime for various odd dimensions is:

{\centering
	\begin{tabular}{|c|c|c|c|c|c|c|c|}
		\hline 
		\multicolumn{8}{|c|}{\textbf{Table 2. Inside the strong coupling regime}} \\ 
		\hline 
		Dimensions & $\Lambda_3\beta$ & $\Lambda_5\beta$ & $\Lambda_7\beta$ & $\Lambda_9\beta$ & $\Lambda_{11}\beta$ & $\Lambda_{13}\beta$ & $m_*\beta$ \\ 
		\hline 
		3 & $0$ & $-$ & $-$ & $-$ & $-$ & $-$ & $0.9624=\ln\phi^2$ \\ 
		\hline 
		5 & $4.1252$ & $0$ & $-$ & $-$ & $-$ & $-$ & $2.03185\approx\ln(3\phi^2)$ \\ 
		\hline 
		7 & $16.3642$ & $6.0403$ & $0$ & $-$ & $-$ & $-$ & $2.89218\approx\ln(7\phi^2)$ \\ 
		\hline
		9 & $65.0883$ & $9.8758$ & $3.8873$ & $0$ & $-$ & $-$ & $3.68896\approx\ln(15\phi^2)$ \\ 
		\hline 
		11 & $...$ & $...$ & $...$ & $...$ & $0$ & $-$ & $4.46019\approx\ln(33\phi^2)$\\
		\hline
		13 & $...$ & $...$ & $...$ & $...$ & $...$ & $0$ & $5.21997\approx\ln(71\phi^2)$\\
		\hline
		$\infty$ & $\infty$ &  $....$ & $....$ & $...$ & $...$ & $...$ & $\infty$ \\
		\hline 
	\end{tabular}\par}
	
where $\phi=1.618$ is the golden ratio and the approximation of the $m_*\beta$ condensate mass at $5d$ is $98.5\%$, at $7d$ is $99.4\%$, at $9d$ is $99.5\%$, at $11d$ is $99.9\%$ and at $13d$ is $99.9\%$...

Interestingly, the fermion condensate mass has a general value of the form:

\begin{equation}
\label{mass}
\boxed
{m_*\beta=\ln(\alpha_n\phi^2)}
\end{equation}
where $\alpha_n=1,3,7,15,33,71,151,319...$, ($n=1,2,3..$ and $\alpha_1=1$). These values are of the form:

\begin{equation}
\alpha_{n+1}=2\alpha_n+\alpha_k
\end{equation}

where $\alpha_k=1,1,1,3,5,9,17,29...$, ($k=1,2,3..$) from \cite{OEIS} with $\frac{\alpha_{k+1}}{\alpha_{k}}=\phi$ and $\alpha_k$ has the growth rate of the Fibonacci numbers.
We see that as dimension increases the $3d$ theory at its strong coupling version always contributes the $\ln\phi^2$ part. All the other dimensions contribute $\ln3, \ln7, \ln15, ...$. By using the changes in the included $\Lambda$ parts in the $5d$, $7d$ and $9d$ gap equations that mentioned before, with the associated mass scales ${\cal M}$, I end up with the corresponding picture:

{\centering
	\begin{tabular}{|c|c|c|c|c|}
		\hline 
		\multicolumn{5}{|c|}{\textbf{Table 3. Gap equations for arbitrary odd dimensions}} \\ 
		\hline 
		Dimensions & $gap\,\,1$ & $gap\,\,2$ & $gap\,\,3$ & $gap\,\,4$ \\ 
		\hline 
		3 & $D_1(-z*)=0$ & $-$ & $-$ & $-$ \\ 
		\hline 
		5 & $D_1(-z*)-0.8754=0$ & $D_3(-z*)=0$ & $-$ & $-$ \\ 
		\hline 
		7 & $D_1(-z*)-1.3890=0$ & $D_3(-z*)+1.5589=0$ & $D_5(-z*)=0$ & $-$ \\ 
		\hline
		9 & $D_1(-z*)-1.8692=0$ & $D_3(-z*)+3.8933=0$ & $D_5(-z*)-2.1527=0$ & $D_7(-z*)=0$ \\ 
		\hline 
		$\infty$ & $....$ &  $....$ & $....$ & $....$ \\
		\hline 
	\end{tabular}\par}
	
We see that as the dimension increases we need larger values of the cut-off (more energy) to regulate the theories. But according to asymptotic freedom of the Gross-Neveu model, higher energies are equivalent to weaker coupling between the fermions. This does not seem to happen when going from $5$ to $7$ dimensions (at the $9d$ theory). Although the cut-off $\Lambda_5\beta$ is larger than $\Lambda_7\beta$ (so the $5$ dimensional theory is in larger energy conditions and weaker coupling) the gap equation that separates the weak from the strong coupling implies that the $5$ dimensional theory is deeper in the strong coupling regime than the $7$ dimensional theory ($3.8933>2.1527$ or ${\cal M}_5^{(9)}\beta^3>{\cal M}_7^{(9)}\beta^5$). It would be interesting to further examine this anomaly in the future, since it might be another mechanism than strongly coupled dynamics that give rise to Fermi condensation like \cite{Chu} or instanton configurations that behave like particles like \cite{Genolini}.	

\section{Free energy at the full bosonisation points $m_*\beta=\ln(\alpha_n\phi^2)$ and $\beta\alpha_*=\pi$}

Let's see for example the free energy density for the model at $5d$ by using a result from \cite{Filothodoros:2018}.

\begin{align}
\label{GNfe5}
\frac{1}{N{\rm Tr}{\mathbb I}_{4}}\Delta f_{(5)} =&-\frac{m_*^2}{2G_{5,*}}+\frac{\beta^4 m_*^4\gamma}{24\pi^3 \beta^5}+\frac{3}{4\pi^2\beta^5}\left[D_5(-z_*)-\frac{1}{24}\ln^4|z_*|D_1(-z_*)\right]\nonumber\\
&+\frac{\beta\alpha*}{8\pi^2\beta^5}\left[D_4(-z_*)+\frac{1}{6}\ln^2|z_*|D_2(-z_*)\right]\,.
\end{align}

The term that appears in the mass gap equation is $\frac{\beta^4 m^4\gamma}{24\pi^3 \beta^5}-\frac{1}{32}\ln^4|z_*|D_1(-z_*)$ and it is zero for $m_*=\ln(3\phi^2)$ and $\beta\alpha_*=\pi$ according to the $5d$ case above. For the system to be in the chirally broken phase at $T=0$ I set as usual the corresponding value analogue to ${\cal M}_5$ (instead of the renormalized $G$ part) which is zero at the critical point. The last term is the charge $Q_5$ where at the full bosonisation point is zero and there is an absence of charged excitations with the canonical partition function of the system to be given by the grand canonical partition function of the same system at fixed imaginary chemical potential \cite{Filothodoros:2016txa}. So the free energy at $5d$ and at the full bosonisation point $\alpha_*=\pi/\beta$ and $m_*\beta=\ln(3\phi^2)$ is 

\begin{equation}
\frac{1}{N{\rm Tr}{\mathbb I}_{4}}\Delta f_{(5)}=\frac{3}{4\pi^2\beta^5}D_5(e^{-\ln(3\phi^2)})\rightarrow \frac{\Delta f_{(5)}}{N}=\frac{3}{{\pi^2\beta^5}}D_5(e^{-\ln(3\phi^2)})\approx 3\frac{4\zeta(5)}{5\pi^2\beta^5}
\end{equation}
with $99\%$ accuracy in the result.

I follow a similar procedure for the higher odd dimensions from results in Appendices A and B of \cite{Filothodoros:2018} and I end up with the table above.

{\centering
	\begin{tabular}{|c|c|c|c|c|c|c|}
		\hline 
		\multicolumn{7}{|c|}{\textbf{Table 4. The free energies}} \\ 
		\hline 
		Dimensions & $\frac{\Delta f_{(3)}}{N}$ & $\frac{\Delta f_{(5)}}{N}$ & $\frac{\Delta f_{(7)}}{N}$ & $\frac{\Delta f_{(9)}}{N}$ & $\frac{\Delta f_{(11)}}{N}$ & $m_*\beta$ \\ 
		\hline 
		3 & $\frac{4 \zeta(3)}{5\pi\beta^3}$ & $-$ & $-$ & $-$ & $-$ & $0.9624=\ln\phi^2$ \\ 
		\hline 
		5 & $-$ & $3\frac{4\zeta(5)}{5\pi^2\beta^5}$ & $-$ & $-$ & $-$ & $2.03185\approx\ln(3\phi^2)$ \\ 
		\hline 
		7 & $-$ & $-$ & $15\frac{4\zeta(7)}{5\pi^3\beta^7}$ & $-$ & $-$ & $2.89218\approx\ln(7\phi^2)$ \\ 
		\hline
		9 & $-$ & $-$ & $-$ & $105\frac{4\zeta(9)}{5\pi^4\beta^9}$ & $-$ & $3.68896\approx\ln(15\phi^2)$ \\ 
		\hline 
		11 & $-$ & $-$ & $-$ & $-$ & $945\frac{4\zeta(11)}{5\pi^5\beta^{11}}$ & $4.46019\approx\ln(33\phi^2)$\\
		\hline 
	\end{tabular}\par}
where 

\begin{equation}
\frac{\Delta f_{(d)}}{N}=0.96166/\pi\beta^3, 2.46186/\pi^2\beta^5, 12.13245/\pi^3\beta^7, 85.42275/\pi^4\beta^9, 774.9189/\pi^5\beta^{11}, ...
\end{equation}
for $d=3,5,7,9,11$ dimensions respectively. We observe the fractional number $\tilde{c}=\frac{4}{5}$ that appears in \cite{Sachdev:1993pr} as a factor in all the free energy densities at the full bosonisation point with an approximation with accuracy over $97.6\%$. The generalized form of the free energy densities at the full bosonisation points is:

\begin{equation}
\label{freeenergy}
\frac{\Delta f_{(2n+1)}}{N}=a_n\frac{\tilde{c}\zeta(2n+1)}{\pi^n\beta^{2n+1}}
\end{equation}
where $n=1,2,3,4,..$ and $a_n=1,3,15,105,..,(2n-1)!!$ , with $a_1=1$.
\section{Summary and Discussion}

The main message of this work is that although there are infrared free terms in the gap equations of the odd $d$ Gross-Neveu model at imaginary chemical potential and finite temperature (large $N$ approximation), we can deal with them by using the higher odd dimensional gap equation as a linear combination of the lower dimensional equations. In the whole process I introduce the new idea of the running cut-off $\Lambda_d$ that varies as a function of energy and describes the connection between the theory's parameters with the energy scale. My calculations have unveiled a generalized form of the mass of the fermionic condensate (\ref{mass}) inside the temperature windows of the model and the values of the free energy densities of the form (\ref{freeenergy}) which are analogous to the rational number $\tilde{c}=\frac{4}{5}$ with a very good approximation.

I believe that my results offer a new window into the physics of bosonisation. As I briefly alluded to in the text there is a kind of anomaly in the gap equation in $5d$ which may be due to another mechanism of fermion condensation than the strongly coupled dynamics or instanton configurations. It would be interesting to examine this identity better in future work.

\section*{Acknowledgements}
I would like to thank A. C. Petkou and N. D. Vlachos for their useful comments and help. 

\begin{thebibliography}{00}
  
\bibitem{Fradkin:1994tt}
E.~H. Fradkin and F.~A. Schaposnik, {\it {The Fermion - boson mapping in
  three-dimensional quantum field theory}}, {\em Phys. Lett.} {\bf B338}
  (1994) 253--258, [{http://arxiv.org/abs/hep-th/9407182}].

\bibitem{Wilczek:1981du}
F.~Wilczek, {\it {Magnetic Flux, Angular Momentum, and Statistics}}, {\em
  Phys. Rev. Lett.} {\bf 48} (1982) 1144--1146.

\bibitem{Polyakov:1988md}
A.~M. Polyakov, {\it {Fermi-Bose Transmutations Induced by Gauge Fields}},
  {\em Mod. Phys. Lett.} {\bf A3} (1988) 325.

\bibitem{Karch:2016sxi}
A.~Karch and D.~Tong, {\it {Particle-Vortex Duality from 3d Bosonization}},
  {\em Phys. Rev.} {\bf X6} (2016), no.~3 031043,
  [{http://arxiv.org/abs/1606.01893}].

\bibitem{Murugan:2016zal}
J.~Murugan and H.~Nastase, {\it {Particle-vortex duality in topological
  insulators and superconductors}}, {\em JHEP} {\bf 05} (2017) 159,
  [{http://arxiv.org/abs/1606.01912}].

\bibitem{Seiberg:2016gmd}
N.~Seiberg, T.~Senthil, C.~Wang, and E.~Witten, {\it {A Duality Web in 2+1
  Dimensions and Condensed Matter Physics}}, {\em Annals Phys.} {\bf 374}
  (2016) 395--433, [{http://arxiv.org/abs/1606.01989}].

\bibitem{Kachru:2016rui}
S.~Kachru, M.~Mulligan, G.~Torroba, and H.~Wang, {\it {Bosonization and Mirror
  Symmetry}}, {\em Phys. Rev.} {\bf D94} (2016), no.~8 085009,
  [{http://arxiv.org/abs/1608.05077}].

\bibitem{Meng:2020}
O. Türker, J. Van den Brink, T. Meng, FS. Nogueira, {\it {Bosonization in $2+1$ dimensions via 
Chern-Simons bosonic particle-vortex duality}}, {\em Phys. Rev.} {\bf D102} (2020), 034506,
[{https://arxiv.org/abs/2004.10789}].


\bibitem{Peskin:1977kp}
M.~E. Peskin, {\it {Mandelstam 't Hooft Duality in Abelian Lattice Models}},
  {\em Annals Phys.} {\bf 113} (1978) 122.

\bibitem{Giombi:2011kc}
S.~Giombi, S.~Minwalla, S.~Prakash, S.~P. Trivedi, S.~R. Wadia, and X.~Yin,
  {\it {Chern-Simons Theory with Vector Fermion Matter}}, {\em Eur. Phys. J.}
  {\bf C72} (2012) 2112, [{http://arxiv.org/abs/1110.4386}].

\bibitem{Aharony:2012ns}
O.~Aharony, S.~Giombi, G.~Gur-Ari, J.~Maldacena, and R.~Yacoby, {\it {The
  Thermal Free Energy in Large N Chern-Simons-Matter Theories}},  {\em JHEP}
  {\bf 03} (2013) 121, [{http://arxiv.org/abs/1211.4843}].
  
\bibitem{Genolini}
Genolini, P.B., Tizzano, L. Instantons, symmetries and anomalies in five dimensions. J. High Energ. Phys. 2021, 188 (2021). https://doi.org/10.1007/JHEP04(2021)188

\bibitem{Filothodoros:2016txa}
E.~G. Filothodoros, A.~C. Petkou, and N.~D. Vlachos, {\it {$3d$ fermion-boson map with imaginary chemical potential}}, {\em Phys. Rev.} {\bf D95} (2017),
  no.~6 065029, [{http://arxiv.org/abs/1608.07795}].

\bibitem{Filothodoros:2018}
E.~G. Filothodoros, Anastasios C. Petkou, Nicholas D. Vlachos,
{\it {The fermion–boson map for large d}},
{\em Nuclear Physics B}
{\bf 941} (2019) Pages 195-224,[{http://arxiv.org/abs/1803.05950}].

\bibitem{Thesis Filothodoros}
E. G. Filothodoros, {\it {The fermion-boson map at imaginary chemical potential in odd dimensions}}, [{http://ikee.lib.auth.gr/record/303052/files/GRI-2019-23684.pdf}].

\bibitem{Barkeshli:2014ida}
M.~Barkeshli and J.~McGreevy, {\it {Continuous transition between fractional
  quantum Hall and superfluid states}}, {\em Phys. Rev.} {\bf B89} (2014),
  no.~23 235116.

\bibitem{Zagier1}
D.~Zagier, {\it The dilogarithm function}, {\em In Frontiers in Number Theory,
  Physics and Geometry II,} (2006, P. Cartier, B. Julia, P. Moussa, P. Vanhove
  (eds.), Springer-Verlag, Berlin-Heidelberg-New York) 3--65.

\bibitem{Filothodoros Lattices}
E. G. Filothodoros, {\it {The fermion boson map for large d and its connection to lattice transformations}}, [{
https://doi.org/10.48550/arXiv.2302.07013}].

\bibitem{Gaiotto}
Gaiotto, D., Kapustin, A., Komargodski, Z. et al. Theta, time reversal and temperature. J. High Energ. Phys. 2017, 91 (2017). https://doi.org/10.1007/JHEP05(2017)091.

\bibitem{Zagier2}
D.~Zagier, {\it {The Bloch-Wigner-Ramakrishnan polylogarithm function}}, {\em Math. Ann. 286 (1990)} (1990), no.~1-3 613--624.

\bibitem{Fosco:1998cq}
C.~D. Fosco, G.~L. Rossini, and F.~A. Schaposnik, {\it {Induced parity breaking term in arbitrary odd dimensions at finite temperature}},  {\em Phys. Rev.} {\bf D59} (1999) 085012, [{http://arxiv.org/abs/hep-th/9810199}{{\tt
hep-th/9810199}}].

\bibitem{Alvarez}
L. Alvarez-Gaume and D. Orlando and S. Reffert, {\it Large charge at large N}, {\em Journal of High Energy Physics} {\bf 12} (2019), [{https://doi.org/10.1007/JHEP12(2019)142}].

\bibitem{Etienne}
R. J. Etienne, {\it {On the Zeros and Extrema of Generalised Clausen Functions }}, {\em Lecture Notes of  TICMI} {\bf vol. 22} (2021), 91-113 ,
  [{http://www.viam.science.tsu.ge/others/ticmi/lnt/vol22/6-R. J. Etienne}].


\bibitem{Petkou:1998wd}
A.~C. Petkou and M.~B. Silva~Neto, {\it {On the free energy of three-dimensional CFTs and polylogarithms}}, {\em Phys. Lett.} {\bf B456} (1999) 147--154, [{http://arxiv.org/abs/hep-th/9812166}].

\bibitem{Christiansen:1999uv}
H.~R. Christiansen, A.~C. Petkou, M.~B. Silva~Neto, and N.~D. Vlachos, {\it {On
  the thermodynamics of the (2+1)-dimensional Gross-Neveu model with complex
  chemical potential}}, {\em Phys. Rev.} {\bf D62} (2000) 025018,
  [{http://arxiv.org/abs/hep-th/9911177}].


\bibitem{Sachdev:1993pr}
S.~Sachdev, {\it {Polylogarithm identities in a conformal field theory in three-dimensions}}, {\em Phys. Lett.} {\bf B309} (1993) 285--288, [{http://arxiv.org/abs/hep-th/9305131}].


\bibitem{DiVecchia}
A.~D'Adda, M.~Luscher, and P.~D. Vecchia, {\it A 1/n expandable series of nonlinear sigma models with instantons}, {\em Nucl. Phys.} {\bf B146} (1978)
  63--76.

\bibitem{Borwein}
J.~M. Borwein and A.~Straub, {\it {Relations for Nielsen polylogarithms}}, {\em J. Approx. Theor.} {\bf 193} (2015) 74--88.

\bibitem{Kolbig}
S.~K. Kolbig, {\it Nielsen's generalized polylogarithms}, {\em SIAM J. Math. Anal.} {\bf 17} (1986) 1232--1258.

\bibitem{ZinnJustin:2002ru}
J.~Zinn-Justin, {\it {Quantum field theory and critical phenomena}},  {\em Int. Ser. Monogr. Phys.} {\bf 113} (2002).

\bibitem{OEIS}
OEIS Foundation Inc. (2023), Entry A298338 in The On-Line Encyclopedia of Integer Sequences, https://oeis.org/A298338.

\bibitem{Chu}
Chu, Chong-Sun and Miao, Rong-Xin, {\it {Fermion condensation induced by the Weyl anomaly}}, {\em PhysRevD.102.046011}, [{https://link.aps.org/doi/10.1103/PhysRevD.102.046011}].


 
\end{thebibliography}
\end{document}